% Notes of the Montreal seminar, 1970 (cote 129, typed part): archive cote 129, site page 59 (PDF page 60).
% typescript. Draft compared: none.
% AI-transcribed by Claude Opus 5.5 (claude-opus-5-5, adaptive thinking, effort high), Message Batches API, batch mtl129_full (msgbatch_01AGnDKHwy6ATCM7ehUa6263), 2026-11-13; not yet checked by eye.
\scanpage{129/59}
\mspage{III-4}

l'épimorphisme induit par la multiplication par $p$. On a
\[ \gamma^{i} = \lambda^{i-1} \circ \lambda^{i-2} \circ \ldots \circ \lambda^{o}, \quad 1 \leq i \leq n-1 \]
on en déduit que l'épimorphisme $\gamma$ est un isomorphisme si et seulement
si $\gamma^{n-1}$ est un monomorphisme, i.e. si on a (e).

(e) $\Rightarrow$ (g) $\Rightarrow$ (c) Avec les notations précédentes, on a:
\[ \beta^{i} = \deleted{J}\inserted{j} \circ \lambda^{n-2} \circ \lambda^{n-3} \circ \ldots \circ \lambda^{i} \] \margin{x \quad j}
(où $j$ désigne l'injection canonique $p^{n-1}G(n) \rightarrow G(1)$). D'après ce qui
précède, (e) entraîne que tous les $\lambda^{i}$ sont des isomorphismes, i.e.
$\beta$ est un monomorphisme. Nous allons montrer par récurrence que si $\beta$ est
un monomorphisme, on a (c), d'où l'on déduit immédiatement que $j$ et donc
$\beta$, sont des isomorphismes.

Si $\beta^{o} = j\circ\gamma^{n-1}$ est un monomorphisme, on a alors $pG(n) = G(n-1)$.
Comme $\beta^{i}$ est un monomorphisme, on a:
\[ [p^{i}G(n)](n-i-1) = p^{i+1}G(n) \ ; \]
mais $p^{i}G(n) = G(n-i)$ par l'hypothèse de récurrence, et donc
$G(n-i-1) = p^{i+1}G(n)$.

2.3. Supposons maintenant que $G(n)$ soit un $\Lambda_{n}$-module plat et qu'il
soit représentable par un schéma fini localement libre sur $S$. Alors $G(i)$,
noyau de la multiplication par $p^{i}$, est représentable, fini et localement
de présentation finie. Mais, $G(n)$ étant supposé $\Lambda_{n}$-plat, on a un épi-%
morphisme $G(n) \overset{p^{n-1}}{\twoheadrightarrow} G(i)$. Fibre à fibre, cette application est surjective

\mspage{59}
